\section{Problem structure and design rationale}
\label{sec:rationale}
\subsection{Three decisions in a budgeted evolutionary search}
A prompt optimizer repeatedly makes three different decisions: which instruction
to generate, which prompt--example outcome to observe, and which candidate to
retain. A good mutation can be lost through an unreliable survivor decision;
a careful survivor rule can be ineffective if the variation operator never
produces useful alternatives. Separating these decisions makes the experimental
question more precise. PACE changes survival while retaining the same operator
family and evaluation policy as its primary control. Any downstream change
includes the resulting feedback through later parents and evaluations.

The evaluation policy determines the observation pattern. If every prompt were
evaluated on all 160 search examples, accuracy estimates would be comparable on
that finite set, but the number of affordable prompts would be small. Conversely,
a large number of lightly evaluated candidates offers more variation while
making rankings less reliable. The budget limits both how many prompts can be explored and how much evidence
is available when survivors are chosen.
PACE operates at that moment. It uses observations already available, rather
than requesting a separate validation batch whenever a survivor decision is due.

A fourth decision occurs at deployment: selecting one instruction from the
retained accuracy--cost trade-off. This decision is held fixed by the utility
in Equation~\ref{eq:winner}. Preserving a diverse retained set and selecting a
single winner serve different purposes. A candidate can be valuable for
maintaining the frontier without maximizing the chosen deployment utility.
Reporting only final accuracy would hide this distinction; reporting only
retention would not answer whether users receive a better prompt.

\subsection{Why common examples matter}
Suppose prompts $p$ and $q$ are tested on the same $n$ independent examples from
a fixed target distribution. For this illustrative sampling model,
\begin{equation}
 \operatorname{Var}(\widehat a_p-\widehat a_q)
 =\frac{\operatorname{Var}(z_p)+\operatorname{Var}(z_q)
       -2\operatorname{Cov}(z_p,z_q)}{n}.
 \label{eq:paired-var}
\end{equation}
The covariance term captures shared example difficulty. If both prompts succeed
on easy examples and fail on difficult ones, comparing them on common examples
can remove part of this nuisance variation. Pairing is not guaranteed to reduce
variance for every pair: the sign and size of the covariance matter. Equation~\ref{eq:paired-var} assumes the simple sampling model above.
The adaptive search does not inherit a coverage guarantee from this calculation.

For binary correctness, only discordant examples contribute a nonzero paired
difference. Let $n_{10}$ count examples on which $p$ is correct and $q$ is wrong,
and $n_{01}$ count the reverse. Then
\begin{equation}
 \widehat\Delta_{pq}=\frac{n_{10}-n_{01}}{n},\qquad
 n=n_{00}+n_{01}+n_{10}+n_{11}.
 \label{eq:discordance}
\end{equation}
This representation distinguishes evidence of equality on the observed examples
from a lack of opportunities to observe disagreement. A pair can have eight
common examples and no discordances; that does not establish equality on future
examples. The bootstrap resamples the observed support and cannot invent an
unseen failure mode.

\subsection{Dominance, preference, and uncertainty}
In the two-objective problem, prompt $p$ point-dominates $q$ when
$\widehat a(p)\geq\widehat a(q)$ and $\widehat c(p)\leq\widehat c(q)$, with at
least one strict inequality. NSGA-II ranks candidates by these estimated
relations and uses crowding to preserve spread within a partially accepted
front~\cite{nsga2}. A prompt with higher accuracy but higher cost may be
non-dominated even when another prompt wins under the fixed deployment utility.
Dominance and preference are therefore not interchangeable criteria.

PACE applies a narrower sufficient condition for evidence-based removal: the
potential dominator must have no greater estimated cost and a strictly positive
lower bound on paired accuracy difference. This rule does not confidence-certify
the entire Pareto ordering. In particular, it does not evidence-eliminate a
cheaper prompt whose paired accuracy advantage is unresolved, nor does it use a
confidence interval for token cost. The retained set can include competitive prompts, poorly measured prompts,
and prompts favored by crowding-based truncation. Retention alone does not
validate a candidate's quality.

\subsection{What the bounded archive is intended to do}
An unresolved prompt has option value only if retaining it can affect a later
decision. It may acquire more observations, re-enter the active population, or
remain available for final winner selection. None of those events is guaranteed.
Archive slots also consume a fixed capacity: protecting one uncertain candidate
can force the eviction of another. This makes the archive a resource allocation
mechanism, not free insurance against noisy selection.

We measure both retention and final performance. Offline reference-Pareto recall measures
whether known finite-pool competitors remain accessible. Online held-out
performance measures the complete search outcome after retention affects
subsequent evolution. A lower false-elimination rate can coexist with unchanged
or worse online accuracy. The experiment is designed to expose that possibility
rather than define success solely through the intermediate retention metric.
